\documentclass[11pt,a4paper]{article}
\usepackage[T1]{fontenc}
\usepackage[utf8]{inputenc}
\usepackage{lmodern}
\usepackage{amsmath,amssymb}
\usepackage[margin=25mm]{geometry}
\usepackage[hidelinks]{hyperref}
\hypersetup{pdftitle={A nonlinear Newton time cell: controlled fluctuations and induced stiffness},pdfauthor={navstokgap research working notes}}
\newcommand{\tightlist}{\setlength{\itemsep}{0pt}\setlength{\parskip}{0pt}}
\setlength{\emergencystretch}{3em}
\setcounter{secnumdepth}{0}
\begin{document}
\section{A nonlinear Newton time cell: controlled fluctuations and
induced
stiffness}\label{a-nonlinear-newton-time-cell-controlled-fluctuations-and-induced-stiffness}

\textbf{Result, 2026-10-27; written derivation, unrefereed.} An
arbitrary intermediate time can be eliminated exactly from a positive
Euclidean mechanical weight. For a convex potential, the normalized
nonlinear conditional law has an explicit fluctuation bound, and the
retained interaction has bounded curvature even when the original force
derivative is unbounded. Both integration and classical minimization
soften the curvature. They differ at a quartic flat centre: positive
supplied \(\hbar\) generates strictly positive quadratic curvature,
while classical elimination retains zero quadratic curvature.

This is a concrete \(0+1\) use of Balaban's instruction to retain the
full normalized interaction and its large-field contribution. It
supplies a nonlinear calculation beyond the Gaussian reference of the
\href{balaban-newton-time-blocking.md}{time-blocking note}. The positive
weight and its action multiplier are premises. An experimentally
inferred positive curvature would need additional arguments before being
identified with a necessary universal quantum action.

\textbf{Targeted check (GPT-6.1 Terra, 2026-10-27).} Terra read this
note and checked the normalized moments/tails, variance and curvature
identities, quartic comparison and curvature constants, and one-halving
conditional independence without finding a mathematical error under the
stated hypotheses. This is a targeted written check; no independent
validation of the physical weighting is claimed.

The Gaussian-convolution and quadratic-infimum constructions are
standard; we prove the cell statements below and claim no literature
novelty. Source motivation:
\href{../docs/balaban/Dimock_BalabanI_1108.1335v2.md}{Dimock I},
sections 2.1, 4.5 and Theorem 14, and
\href{../docs/balaban/Dimock_BalabanII_1212.5562v2.md}{Dimock II},
Theorem 3.1 and sections 3.13--3.14 (selected passages read). The old
\href{rivero-1998-conjecture-central-forces.md}{quartic limit theorem},
section 3, uses \(V=\lambda q^4\); here the convention is
\(V=\lambda q^4/4\). The two normalizations must not be mixed.

\subsection{1. The complete one-vertex
weight}\label{the-complete-one-vertex-weight}

Let \(m,u,v,\hbar>0\), \(T=u+v\), and fix positions \(x,y\). Put

\[
 \mu=\frac{vx+uy}{T},\qquad
 a=\frac{mT}{uv},\qquad w=\frac T2.
 \tag{1}
\]

Here \(a\) has units of action per length squared and \(w\) is a time.
Completing the kinetic square gives

\[
 \frac{m(z-x)^2}{2u}+\frac{m(y-z)^2}{2v}
 =\frac{m(y-x)^2}{2T}+\frac a2(z-\mu)^2.
 \tag{2}
\]

The symmetric Euclidean action also includes \(uV(x)/2+wV(z)+vV(y)/2\).
Define

\[
 I_\hbar(\mu)=\sqrt{\frac{a}{2\pi\hbar}}
 \int_{\mathbb R}\exp\left[-\frac{F_\mu(z)}{\hbar}\right]dz,
 \quad F_\mu(z)=\frac a2(z-\mu)^2+wV(z),\quad
 A_\hbar=-\hbar\log I_\hbar.
 \tag{3}
\]

Then the integrated two-cell kernel, with the free normalization at each
cell, is exactly

\[
 \sqrt{\frac{m}{2\pi\hbar T}}
 \exp\left[-\frac1\hbar\left\{
 \frac{m(y-x)^2}{2T}+\frac u2V(x)+\frac v2V(y)
       +A_\hbar(\mu)\right\}\right].
 \tag{4}
\]

We retain the whole function \(A_\hbar\); it is not generally a quartic
polynomial. The normalized conditional law \(\nu_{\hbar,\mu}\) of \(Z\)
is \(e^{-F_\mu/\hbar}dz\) divided by its integral. No unnormalized
action cost is used as a probability estimate.

For comparison define

\[A_0(\mu)=\min_z F_\mu(z).\tag{5}\]

This is a classical Euclidean stationary elimination. It is not the
real-time Newton action: for real time the potential sign in (2)--(4)
changes and an oscillatory integral replaces the positive weight. In
particular, these probability bounds cannot simply be continued to real
time. No \(c\) or relativistic premise is used here.

\subsection{2. Convexity gives normalized
control}\label{convexity-gives-normalized-control}

Assume \(V\in C^2(\mathbb R)\) is convex, with sufficient growth
regularity for the following integrations by parts. Polynomials with
nonnegative second derivative satisfy the hypotheses. The Gaussian
kinetic term makes \(F_\mu\) strongly convex, so it has one minimizer
\(z_*(\mu)\) and \(F_\mu''\ge a\). All moments needed below are finite.

\textbf{Proposition 1.} For \(n=1,2,\ldots\) and \(r>0\),

\[
 \mathbb E_{\nu_{\hbar,\mu}}(Z-z_*)^{2n}
       \le(2n-1)!!\left(\frac\hbar a\right)^n,
 \qquad
 \nu_{\hbar,\mu}(|Z-z_*|\ge r)
       \le2\exp\left[-\frac{ar^2}{2\hbar}\right].
 \tag{6}
\]

Also

\[
 \operatorname{Var}(Z)\le\frac\hbar a,\qquad
 A_\hbar'=a(\mu-\mathbb EZ),\qquad
 A_\hbar''=a-\frac{a^2}{\hbar}\operatorname{Var}(Z),\qquad
 0\le A_\hbar''\le a.
 \tag{7}
\]

\emph{Proof.} Write \(Y=Z-z_*\). Strong convexity and \(F_\mu'(z_*)=0\)
imply \(YF_\mu'(z_*+Y)\ge aY^2\). Integration by parts gives

\[
 \mathbb E[Y^{2n-1}F_\mu'(Z)]
       =(2n-1)\hbar\mathbb EY^{2n-2}.
 \tag{8}
\]

The left side is at least \(a\mathbb EY^{2n}\). Induction proves the
moment bound. Expanding \(\cosh(tY)\) in its nonnegative even powers
gives \(\mathbb E\cosh(tY)\le e^{\hbar t^2/(2a)}\). Since
\(e^{t|Y|}\le e^{tY}+e^{-tY}\), exponential Markov inequality with
\(t=ar/\hbar\) gives the tail bound.

For variance, take independent copies \(Z,Z'\). Integration by parts
also gives \(\mathbb EF_\mu'=0\) and
\(\operatorname{Cov}(Z,F_\mu'(Z))=\hbar\). Consequently

\[
 \hbar=\frac12\mathbb E[(Z-Z')(F_\mu'(Z)-F_\mu'(Z'))]
       \ge a\operatorname{Var}(Z).
 \tag{9}
\]

Differentiating the normalized integral (3) gives the two derivative
identities in (7); in particular
\(\partial_\mu\mathbb EZ=a\operatorname{Var}(Z)/\hbar\). This proves the
curvature bounds. \(\square\)

If additionally \(0\le V''\le K\), the same two-copy identity with the
upper slope bound gives

\[
 \frac\hbar{a+wK}\le\operatorname{Var}(Z)\le\frac\hbar a,
 \qquad
 0\le A_\hbar''\le\frac{awK}{a+wK}.
 \tag{10}
\]

These are bounds for the full normalized nonlinear conditional law. They
require neither a quadratic approximation nor a small quartic coupling.

The function \(A_\hbar/w\), viewed as a potential of the coordinate
\(\mu\), has force derivative bounded by \(a/(mw)=2/(uv)\). This is a
mathematically well-defined coarse potential. Identifying its gradient
with the force of a physical Newtonian body is a further modeling
choice: (4) is an endpoint kernel, not an autonomous Newton equation
derived from it.

\subsection{3. Classical elimination has the same curvature
bound}\label{classical-elimination-has-the-same-curvature-bound}

The minimizer satisfies

\[a(z_*-\mu)+wV'(z_*)=0,\qquad
 z_*'=\frac a{a+wV''(z_*)}.
 \tag{11}
\]

The implicit function theorem and direct differentiation yield

\[
 A_0'=a(\mu-z_*),\qquad
 A_0''=\frac{awV''(z_*)}{a+wV''(z_*)},\qquad
 0\le A_0''\le a.
 \tag{12}
\]

Thus bounded coarse curvature is already possible in classical
elimination. For the convex quartic \(V=\lambda z^4/4\), \(\lambda>0\),
the positive-\(\hbar\) interaction approaches this classical envelope
with an explicit bound:

\[
 0\le A_\hbar-A_0
 \le\frac{3w\lambda z_*^2\hbar}{2a}
        +\frac{3w\lambda\hbar^2}{4a^2}.
 \tag{13}
\]

\emph{Proof.} After centering the integral at \(z_*\), let
\(Y\sim N(0,\hbar/a)\) and \(R_V=V(z_*+Y)-V(z_*)-V'(z_*)Y\ge0\).
Cancellation of the linear terms in \(F_\mu\) gives

\[
 A_\hbar=A_0-\hbar\log\mathbb E
                   e^{-wR_V/\hbar}.
 \tag{14}
\]

The expectation is at most one, and Jensen's inequality bounds the
additional action by \(w\mathbb ER_V\). For the quartic, the Gaussian
second and fourth moments give exactly the right side of (13). This
proves local uniform convergence as \(\hbar\downarrow0\). \(\square\)

Classical minimization and positive integration therefore share one
regularity mechanism. A criterion requiring only a finite ODE Lipschitz
constant would not distinguish them in this convex cell.

\subsection{4. The quartic flat centre does distinguish the two
weights}\label{the-quartic-flat-centre-does-distinguish-the-two-weights}

For \(V=\lambda z^4/4\) and \(\mu=0\), symmetry gives \(\mathbb EZ=0\)
and \(z_*=0\). Write

\[g=\frac{w\lambda\hbar}{a^2},\qquad
 \chi_\hbar=A_\hbar''(0).\tag{15}\]

The classical curvature is \(A_0''(0)=0\). For every supplied
\(\hbar>0\), however,

\[
 0<a\frac{\sqrt{1+4g}-1}{\sqrt{1+4g}+1}
       \le\chi_\hbar\le a\min(1,3g).
 \tag{16}
\]

Moreover \(\chi_\hbar/(ag)\to3\) as \(g\downarrow0\).

\emph{Proof.} Integration of \((ze^{-F_0/\hbar})'\) gives

\[a\mathbb EZ^2+w\lambda\mathbb EZ^4=\hbar,\qquad
 \chi_\hbar=\frac{aw\lambda}{\hbar}\mathbb EZ^4.\tag{17}\]

Let \(r=a\mathbb EZ^2/\hbar\). Jensen's inequality
\(\mathbb EZ^4\ge(\mathbb EZ^2)^2\) gives \(r+gr^2\le1\), hence
\(r\le2/(1+\sqrt{1+4g})\). Equation (7) gives \(\chi_\hbar=a(1-r)\),
proving the lower bound. Equation (6) at \(n=2\) gives
\(\mathbb EZ^4\le3(\hbar/a)^2\), proving the upper bound together with
\(\chi_\hbar\le a\).

Under \(Z=\sqrt{\hbar/a}\,X\), the normalized density of \(X\) is
proportional to \(e^{-X^2/2-gX^4/4}\). Its fourth moment tends to the
Gaussian value 3 as \(g\downarrow0\) by dominated convergence. Equation
(17) gives the stated limit. \(\square\)

This is a finite nonlinear example in which integrating fluctuations
generates a restoring quadratic term at a classically flat point. The
local frequency associated with the defined coarse potential is
\(\omega_{\mathrm{cell}}^2=\chi_\hbar/(mw)\), so

\[
 \omega_{\mathrm{cell}}^2
       \sim\frac{3\lambda\hbar}{ma}\quad(g\downarrow0).
 \tag{18}
\]

The word frequency here describes this coarse-potential diagnostic; it
is not an asserted spectral gap or a real-time oscillation frequency
without a reconstruction argument.

A direct insertion calculation verifies the generated term. At fixed
\(a,w,\hbar,\mu\), the right derivative at \(\lambda=0\) is

\[
 \left.\partial_\lambda A_\hbar(\mu)\right|_{0+}
 =\frac w4\left[\mu^4+6\mu^2\frac\hbar a
                         +3\left(\frac\hbar a\right)^2\right].
 \tag{19}
\]

This follows by differentiating (3) under the Gaussian integral. It is a
statement about a derivative, not a claim that the full perturbation
series converges. It agrees with the coefficient 3 in (18).

\subsection{5. One nonlinear halving and its large-field
events}\label{one-nonlinear-halving-and-its-large-field-events}

In a finite nearest-neighbour Euclidean time chain with fixed endpoint
positions, use the normalized joint law and choose a set \(\mathcal I\)
of interior vertices with no two adjacent, retaining all other vertices.
Conditioned on the retained values, the variables \(Z_j\),
\(j\in\mathcal I\), are independent, each with a density of the form
(3), with its own \(a_j,w_j,\mu_j\). The exact retained kernel remains a
product of nearest-neighbour kernels, although their actions need not
have the original kinetic-plus-potential form.

Define a bad event relative to each cell's conditional minimizing value:

\[
 B_j=\left\{|Z_j-z_{*,j}|>b\sqrt{\hbar/a_j}\right\},\qquad
 p_b=2e^{-b^2/2}.
 \tag{20}
\]

For any \(H\subset\mathcal I\), Proposition 1 and conditional
independence give, including after averaging the retained variables,

\[\mathbb P\left(\bigcap_{j\in H}B_j\right)\le p_b^{|H|}.
 \tag{21}\]

For \(b>\sqrt{2\log2}\) this is a decaying bound. It controls the
normalized nonlinear measure, including configurations outside the
chosen tubes. It is one halving statement for this convex local action.
After further halvings the generated interactions must be controlled; we
have not proved that the same original-potential formula closes.

\subsection{6. Written controls and physical
meaning}\label{written-controls-and-physical-meaning}

For \(V=0\), (3) gives \(A_\hbar=0\) and variance \(\hbar/a\),
saturating (7). For \(V=kz^2/2\), \(k\ge0\), completing the square gives

\[
 A_\hbar(\mu)=\frac{awk}{2(a+wk)}\mu^2
                 +\frac\hbar2\log\left(1+\frac{wk}{a}\right).
 \tag{22}
\]

Its curvature agrees exactly with (10) and (12); the fluctuation term is
independent of \(\mu\) and leaves that curvature unchanged. This
quadratic control distinguishes a force-curvature effect from the
quartic fluctuation-generated curvature in section 4.

For \(V=-Fz\), (3) gives

\[
 A_\hbar=-wF\mu-\frac{w^2F^2}{2a},\qquad
 \frac{w^2F^2}{2a}=\frac{F^2uvT}{8m}.
 \tag{23}
\]

This recovers the arbitrary-cut Galileo correction in the time-blocking
note, with its normalization explicit.

The cell curvature in (16) is a possible scale diagnostic, conditional
on the entire weighting model. If \(a,w,\lambda\) are independently
fixed and \(0<\chi<a\) is the measured value of the defined cell
curvature, inverting (16) gives the model-conditional action interval

\[
 \frac{a\chi}{3w\lambda}\le\hbar
       \le\frac{a\chi}{w\lambda(1-\chi/a)^2}.
 \tag{24}
\]

Indeed, with \(q=\chi/a\), (16) gives \(q/3\le g\) and
\(\sqrt{1+4g}\le(1+q)/(1-q)\), so \(g\le q/(1-q)^2\). The lower endpoint
is asymptotically sharp in the small-\(g\) limit. This inference
presupposes that the measured response is the response of (3), rather
than an unmodeled apparatus effect.

A nonzero measured value could also come from a classical statistical
weight with an independently supplied action-like noise parameter. ODE
regularity or this positive Euclidean integral alone cannot identify
that parameter with the real-time phase constant, a universal apparatus
restriction, or Planck's measured value. Those are separate physical
tests. The bounded full-cell result is useful because it makes the
alternatives quantitatively testable.

\subsection{7. Consequence for STATE}\label{consequence-for-state}

The first nonlinear time-elimination step is now controlled without
replacing its law by a Gaussian: normalized tails, endpoint response,
coarse curvature and a simultaneous one-halving bad-event bound are
explicit. The convex quartic generates positive flat-centre curvature at
supplied \(\hbar>0\), whereas classical elimination remains regular with
zero flat-centre curvature. Next investigate branch competition or
real-time/record reconstruction of this diagnostic. The Newton necessity
question remains active; no positivity of \(\hbar\) is derived from the
choice of a positive weight. These are written, unrefereed calculations.
\end{document}
